\documentclass{homework}
% \usepackage{graphicx} % Required for inserting images
\usepackage{amsmath}
\usepackage{gensymb}
\usepackage{subfig}
\usepackage{float}

\class{ECEN 5448}
\title{Homework 5}
\author{Sean Jennings}
\date{December 19, 2023}
\newcommand{\control}{\mathcal{C}}

\begin{document}
\maketitle

\section{Problem 1}
\subsection{}
We note that the system is already in the Kalman Decomposition form (KDF):
\begin{equation}
    \dot{x} = \mat{
        A_{11} & A_{12} \\
        0 & A_{22}
    } x + \mat{
        B_1 \\ 0
    } u
    \label{eq:kdf}
\end{equation}
with
\begin{align*}
    A_{11} &= \mat{0 & 1 \\ 0 & 0} \qquad A_{12} = \mat{2 \\ 3} \\
    A_{22} &= \mat{-1} \qquad B_1 = \mat{0 \\ 1}.
\end{align*}
$A_{22}$ is Hurwitz so the uncontrollable subsystem is exponentially stable
and the system is therefore feedback stabilizable.

Furthermore, we note that the controllable subsystem is already in controllability
canonical form (CCF).

To find $K\in\R^{1\times 3}$, we first decompose it into $K = \mat{K_1 K_2}$, with
$K_1 \in \R^{1 \times 2}$, $K_2\in\R$ so that the closed-loop system with $u=-Kx$
is given by:
\begin{equation*}
    \dot{x} = \mat{
        A_{11} - B_1 K_1 & A_{12} - B_1 K_2 \\
        0 & A_{22}
    } x
\end{equation*}
Since $K_2$ does not affect the eigenvalues of our closed loop system, we arbitrarily
set it to 0 and will select $K_1 = [k_0, k_1]$ ($k_0,k_1\in\R$) such that the poles of 
the controllable subsystem are $\lambda = \lambda_{1,2} = -1$. Thus, the desired characteristic polynomial
is given by $(\lambda + 1)^2 = \lambda^2 + 2\lambda + 1$. The characteristic polynomial
of matrix $A_{11} - B_1K_1$ is given by:
\begin{equation*}
    \det(\lambda I - (A_{11} - B_1K_1)) = \det{\lambda I - \mat{0 & 1 \\ -k_0 & -k_1}} = \lambda^2 + k_1 \lambda + k_0
\end{equation*}

Therefore, we select $k_0 = 1$, $k_1=2$ to achieve the desired poles, yielding a final result of
\begin{equation*}
    K = \mat{1 & 2 & 0}
\end{equation*}
so that the feedback-stabilized closed-loop system has eigenvalues $\lambda_{1,2,3} = -1$.

We verify in Python that all eigenvalues of the matrix $A - BK$ are -1.

\subsection{}
\begin{letterenum}
\item The controllability matrix of the system $(A, B)$ is given by:
\begin{equation*}
    \control(A, B) = \mat{
        1 & -1 & 2 \\
        -1 & 2 & 0 \\
        1 & -1 & 0
    }
\end{equation*}
which has determinant det$(\control(A, B)) = -2$, so the system is controllable and
there exists an invertible $P\in\R^{3\times 3}$ such that the system $(A', B') = (PAP^{-1}, PB)$
is in CCF.

In python, we determine the characteristic polynomial is:
\begin{equation*}
    \lambda^3 + \lambda^2 - 2
\end{equation*}
which implies that the CCF system is:
\begin{equation*}
    A' = \mat{
        0 & 1 & 0 \\
        0 & 0 & 1 \\
        2 & 0 & -1
    } \qquad B' = \mat{0 \\ 0 \\ 1}
\end{equation*}
with $P$ given by:
\begin{equation*}
    P = \control(A', B') \control(A, B)^{-1} = \frac{1}{2}\mat{
        1 & 0 & -1 \\
        -1 & 2 & 3 \\
        1 & 0 & 1
    }.
\end{equation*}
\item 
Let $z = Px$ be the change of coordinates into CCF. In terms of the change of coordinates,
we have $u = -Kx = -KP^{-1}z$. Then, the closed-loop
dynamics of $(A', B')$ are given by:
\begin{equation*}
    \dot{z} = A'z - B'u = A'z - B'KP^{-1}z = (A' - B'KP^{-1})z.
\end{equation*}
We select $K = \mat{k_0 & k_1 & k_2}P$ with $k_0,k_1,k_2\in\R$ so that the characteristic polynomial (c.poly.)
is given by:
\begin{equation*}
    \det(A' - B'\mat{k_0 & k_1 & k_2} (PP^{-1})) = \det\biggl(\mat{
        0 & 1 & 0 \\
        0 & 0 & 1 \\
        2 - k_0 & -k_1 & -1 - k_2
    }\biggr) = \lambda^3 + (1 + k_2) \lambda ^2 + k_1 \lambda + (k_0 - 2)
\end{equation*}
We place the poles so that the c.poly. is given by:
\begin{equation*}
    (\lambda + 1)(\lambda + (2 - i))(\lambda + (2 + i)) = \lambda^3 + 5\lambda^2 + 9\lambda + 5
\end{equation*}
Equating the polynomials, the gain matrix, $K'\in\R^{1\times 3}$, in the change of coordinate system is:
\begin{equation*}
    K' := \mat{k_0 & k_1 & k_2} = \mat{7 & 9 & 4}
\end{equation*}
\item 
Converting back to the original coordinate system, we have
\begin{equation*}
    K = K'P = \mat{1 & 9 & 12}
\end{equation*}
\end{letterenum}


\subsection{}
Define $x\in\R^3$ such that $x:=\mat{x_1 & x_2 & x_3}^T$ and the function $f: \R^3 \times \R \to \R^3$
to be the given system dynamics, such that $\dot{x} = f(x,u)$. Clearly, $f(0_3, 0) = 0$ so that 
$x^* = 0_3$ and $u^*=0$ is an equilibrium point for the system. To see that the system is unstable with $u = 0$, we check
the eigenvalues of the linearized system with $\tilde{x} := x - x^* = x$, $\tilde{u} = u - u^* = u$:
\begin{align*}
    \dot{\tilde{x}} &= \dot{x} \approx \nabla_x f(x^*, u^*) \tilde{x} + \nabla_u f(x^*, u^*) \tilde{u} \\
        &= \mat{
            -1 & 2x_2 & 0 \\
            0 & -1 + x_3 & x_2 \\
            x_2 & x_1 & 1
        } \eval_{x_1=x_2=x_3=0} \tilde{x} + \mat{1 \\ -1 \\ 1} \tilde{u} \\
        &= \mat{
            -1 & 0 & 0 \\
            0 & -1 & 0 \\
            0 & 0 & 1
        } \tilde{x} + \mat{1 \\ -1 \\ 1} \tilde{u} \\
        &=: A\tilde{x} + B \tilde{u}
\end{align*}
where we have defined the matrix $A\in\R^{3\times 3}$ and $B\in\R^3$ as the dynamics matrix and control matrix of the linearized system,
respectively.  The eigenvalues of $A$ are $\lambda_{1,2} = -1$ and $\lambda_{3} = 1$. Since $A$ has a positive
eigenvalue, by Lyapunov's first theorem, the equilibrium point $x^*$ is locally unstable with $u=0$.

To stabilize the system, we look to convert it to CCF. However, the controllability matrix, given by:
\begin{equation*}
    \control(A, B) = \mat{
        0 & 0 & 0 \\
        1 & -1 & 1 \\
        1 & 1 & 1
    }
\end{equation*}
is singular and we must first find the $P_1\in\R^{3\times 3}$ s.t. the system given by $y=P_1\tilde{x}$ is in KDF.

We select $P_1$ such that:
\begin{align*}
    P_1 v_1 := P_1 \mat{0 \\ 1 \\ 0} = e_1 = \mat{1 \\ 0 \\ 0} \\
    P_1 v_2 := P_1 \mat{0 \\ 0 \\ 1} = e_2 = \mat{0 \\ 1 \\ 0} \\
    P_1 v_3 := P_1 \mat{1 \\ 0 \\ 0} = e_3 = \mat{0 \\ 0 \\ 1}
\end{align*}
where we have defined $v_1,v_2\in\R^3$ such that $\{v_1, v_2\}$ forms a basis for
$\control(A, B)$ and
$v_3\in\R^3$  such that $\{v_1,v_2,v_3\}$ forms a basis for $\R^3$. It is easy to see that
\begin{equation*}
    P_1 = \mat{
        0 & 1 & 0 \\
        0 & 0 & 1 \\
        1 & 0 & 0
    }
\end{equation*}
satisfies the above equations. Our KDF system $(A_{11}, A_{12}, A_{22}, B_1)$
(as defined by eq. \ref{eq:kdf})
has
\begin{equation*}
    A_{11} = \mat{
        -1 & 0 \\
        0 & 1
    } \qquad A_{22} = \mat{-1} \qquad B_1 = \mat{1\\1}
\end{equation*}
$A_{22}$ is again Hurwitz so the system is feedback stabilizable.

Let intervible $P_2\in\R^{2\times2}$ be the change coordinate matrix which brings
the system $(A_{11}, B_1)$ to the CCF form $(A', B')$.
The c.poly. of $A_{11}$ is
$(\lambda + 1)(\lambda - 1) = \lambda^2 - 1$, so $A'$ is given by:
\begin{equation*}
    A' = \mat{
        0 & 1 \\
        1 & 0
    }.
\end{equation*}
$P_2$ is:
\begin{equation*}
    P_2 = \control(A', B') \control(A_{11}, B_1)^{-1} = \frac{1}{2} \mat{
        -1 & 1 \\
        1 & 1
    }
\end{equation*}
and $B' = \mat{0 & 1}^T$.

Now, we have coordinate shifts $P_1$ and $P_2$ which convert the original system
to KDF and then the controllable KDF subsystem to CCF. We next design a feedback
gain for the CCF system and shift it back to our original coordinate system.

Let $z = P_2 y_1$ be the state vector of the CCF system, where $y_1\in\R^2$ contains
the controllable states of $y$. Let $K_1,K'\in\R^{1\times 2}$ be the gain
matrices of the KDF subsystem $(A_{11}, B_1)$ and the CCF system $(A', B')$, respectively. 
We relate $K_1$ and $K'$ to the gain matrix $K\in\R^{1 \times 3}$ of our original system
by $K = \mat{K_1 & 0} P_1$ and $K' = K_1 P_2$, so that our feedback $u=Kx$ converted
to CCF is:
\begin{align*}
    u = Kx &= KP_1^{-1}y \\
        &= \mat{K_1 & 0} y \\
        &= \mat{K_1 P_2^{-1} z & 0} \\
        &= \mat{K' z & 0}.
\end{align*}
Then, with $K' = \mat{k_0 & k_1}$, our CCF system has closed-loop dynamics:
\begin{align*}
    \dot{z} &= (A' + B'\mat{k_0 & k_1}) z \\
        &= \mat{
            0 & 1 \\
            k_0 + 1 & k_1
        }
\end{align*}
which has c.poly. $\lambda^2 - k_1 \lambda - (k_0 + 1)$. Selecting eigenvalues
$\lambda_{1,2} = -1$ yields a desired c.poly. of $(\lambda + 1)^2 = \lambda^2 + 2\lambda + 1$.
Thus, we have
\begin{align*}
    k_0 + 1 = -1 \Rightarrow k_0 = -2 \\
    k_1 = -2,
\end{align*}
and we can now apply change of coordinate matrices to compute $K_1$ and $K$:
\begin{align*}
    K_1 &= \mat{-2 & -2} P_2 = \mat{0 & -2} \\
    K &= \mat{K_1 & 0} P_1 = \mat{0 & -2 & 0} P_1 \\
      &= \mat{0 & 0 & -2}.
\end{align*}
It is easy enough to see that the matrix $A + BK$ given by
\begin{equation*}
    A + BK = \mat{
        -1 & 0 & 0 \\
        0 & -1 & 0 \\
        0 & 0 & 1
    } + \mat{
        0 \\ 1 \\ 1
    } \mat{0 & 0 & -2}
    = \mat{
        -1 & 0 & 0 \\
        0 & -1 & -2 \\
        0 & 0 & -1
    },
\end{equation*}
has the desired poles $\lambda_{1,2,3} = -1$ (since the eigenvalues of a triangular matrix are its diagonal entries).

Thus, the linearized system $A, B$ (about the origin) is globally exponentially stable which implies that the nonlinear system $f$
is asymptotically stable around the origin.

\section{Problem 2}
\subsection{}
To design a Luenberger observer, we must find an $L\in\R^3$ such that the matrix $A + LC$ has the desired eigenvalues.
But this is equivalent to finding a feedback matrix $K\in\R^{1\times 3}$
for the dual system $(A^T, C^T)$ with closed-loop dynamics $A^T + C^T K$, when
we equate $K=L^T$ (since for any square matrix $H$, $H$ and $H^T$ have the same eigenvalues).

We therefore follow the same steps as to find a feedback matrix $K$ for dual
system $(A^T, C^T)$ with:
\begin{equation*}
    A^T = \mat{
        0 & 0 & 1 \\
        1 & 0 & 0 \\
        0 & 1 & 0
    } \qquad C^T = \mat{1 \\ 0 \\ 0}.
\end{equation*}
Let $(A', C')$ be the CCF system, given by:
\begin{equation*}
    A' = \mat{
        0 & 1 & 0 \\
        0 & 0 & 1 \\
        1 & 0 & 0
    } \qquad C' = \mat{0 \\ 0 \\ 1}.
\end{equation*}
and
\begin{equation*}
    P = \control(A', C')\control(A^T, C^T)^{-1} = \mat{
        0 & 0 & 1 \\
        0 & 1 & 0 \\
        1 & 0 & 0
    }
\end{equation*}
Let $K = \mat{k_0 & k_1 & k_2} P$ so that the closed-loop system has c.poly:
\begin{equation*}
    \det(\lambda I - (A' + C' \mat{k_0 & k_1 & k_2})) = \mat{
        0 & 1 & 0 \\
        0 & 0 & 1 \\
        k_0 + 1 & k_1 & k_2 
    } = \lambda^3 - k_2 \lambda^2 - k_1 \lambda - (k_0 + 1)
\end{equation*}
The desired c.poly. is $\lambda^3 + 5\lambda^2 + 9\lambda + 5$. Solving for
$k_0$, $k_1$, and $k_2$ and equating to the observer gain $L$, we have:
\begin{equation*}
    L = K^T = P^T \mat{k_0 \\ k_1 \\ k_2} = \mat{
        0 & 0 & 1 \\
        0 & 1 & 0 \\
        1 & 0 & 0
    } \mat{-6 \\ -9 \\ -5} = \mat{-5 \\ -9 \\ -6}.
\end{equation*}
The observer is given by the system:
\begin{align*}
    \dot{\hat{x}} = A\hat{x} + Bu - L (y - \hat{y}) \\
    \hat{y} = C \hat{x}.
\end{align*}

\subsection{}
\begin{letterenum}
\item We divide the pendulum equation by $m\ell^2$, define $x\in\R^2$ to be
$x = \mat{\theta & \dot{\theta}}^T$, and define $f: \R^2 \times \R \to \R^2$ to be
the system dynamics given by:
\begin{equation*}
    \dot{x} = f(\mat{\dot{\theta} \\ \ddot{\theta}}, u) := \mat{\dot{\theta} \\ -10 \sin\theta} + \mat{0 \\ 1} u
\end{equation*}
The jacobian w.r.t. $x$ and $u$ are given by:
\begin{equation*}
    \nabla_x f(x, u) = \mat{0 & 1 \\ -10\cos\theta & 0} \qquad \nabla_u f(x, u) = \mat{0 \\ 1}
\end{equation*}
so that linearized system ($\tilde{x} := x - x^*$, $\tilde{u} = u - u^*$) about the equilibrium
$x^* := \mat{\pi & 0}^T, u^* := 0$ is given by:
\begin{align*}
    \dot{\tilde{x}} &\approx \nabla_x f(x^*, u^*) \tilde{x} + \nabla_u f(x^*, u^*) \tilde{u} \\
        &= A \tilde{x} + B \tilde{u}
\end{align*}
with observation $\tilde{y}$ given by:
\begin{equation*}
    \tilde{y} := y - y^* = Cx - Cx^* = C\tilde{x}
\end{equation*}
where
\begin{align*}
    A = \mat{
            0 & 1 \\
            10 & 0
        } \qquad B = \mat{0 \\ 1} \qquad C = \mat{1 & 0}.
\end{align*}

By the separation theorem, we can begin by designing an observer with gain $L\in\R^{2\times 2}$
such that the dual system $(A^T, C^T)$ is feedback stable with poles $\lambda_{1,2} = -2$ (
    c.poly of $\lambda^2 + 4 \lambda + 4)$.

Converting to CCF system $(A', C')$, we find $L$:
\begin{align*}
    A' :&= PA^TP^{-1} = \mat{0 & 1 \\ 10 & 0} \\
    C' :&= PC^T = \mat{0 \\ 1} \\
    P :& = \mat{0 & 1 \\ 1 & 0} \\
    L^T :&= \mat{\ell_0 & \ell_1} P \\
    \det(I\lambda - (A' + C' \mat{\ell_0 \ell_1})) &= \lambda^2 - \ell_1 - (10 - \ell_0) \\
    \ell_0 &= -14 \\
    \ell_1 &= -4 \\
    L &= P^T \mat{\ell_0 \\ \ell_1} = \mat{\ell_1 \\ \ell_0} = \mat{-4 \\ -14}.
\end{align*}
Since the primal system $(A, B)$ is equal to the CCF-shifted dual system $(A', C')$
we can take $K = \mat{k_0 & k_1}$ with $k_0 = \ell_0$ and $k_1 = \ell_1$ so that
\begin{equation*}
    K = \mat{-14 & -4}
\end{equation*}

\item The closed loop system of the linearized system (writing $\tilde{x} = x$ and $\tilde{u} = u$ for brevity)
w/ both observer and feedback is given by:
\begin{align*}
    \dot{x} = Ax + Bu \\
    y = Cx \\
    \dot{\hat{x}} = A\hat{x} + Bu - L(y - \hat{y}) \\
    \hat{y} = C\hat{x}
    u = K\hat{x}
\end{align*}
We define the error $e\in\R^2$ as $e:=x - \hat{x}$, which has dyanmics given by:
\begin{align*}
    \dot{e} &= \dot{x} - \dot{\hat{x}} \\
        &= Ax + Bu - (A\hat{x} + Bu - LC(x - \hat{x})) \\
        &= (A + LC) e.
\end{align*}
The state dynamics w.r.t. to the error is given by:
\begin{equation*}
    \dot{x} = Ax + Bu = Ax + BK \hat{x} = Ax + BK(x + e).
\end{equation*}
Stacking state and estimation error results in closed loop system dynamics of:
\begin{equation*}
    \mat{
        \dot{x} \\ \dot{e}
    } = \mat{
        A + BK & BK \\
        0 & A+LC
    } \mat{x \\ e}
    =: M\mat{x \\ e}.
\end{equation*}

So the linearized closed-loop dynamics has eigenvalues given by:
\begin{equation*}
    \text{eig}(M) = \text{eig}(A + BK) \cup \text{eig}(A + LC).
\end{equation*}

Since we designed $A + BK$ and $A + LC$ to each have both eigenvalues equal to -2,
the matrix $M$ is Hurwitz, and the linearized system is globally exponentially stable.
Therefore by Lyapunov's first method, the
equilibrium point $x^* = \mat{\pi & 0}^T$ is locally asymptotically stable.
\end{letterenum}










\end{document}
